% ED26-BM-RAT2021-14-1: exact MIT 2021 chapter, printed 747--755.
\chapter{Scientific Rationality and Objectivity}
\label{sec:p11-rationality-14-1}

Line Edslev Andersen and Hanne Andersen's chapter asks what makes
scientific inquiry rational and objective, and whether the relevant
assessment concerns individual scientists or scientific communities.
The shift in the assessed agent matters: people can disagree while
contributing to a community's inquiry, and a community can agree without
thereby having removed its biases. The source reconstructs several
twentieth-century accounts and applies their differences to peer review;
it does not announce a single method that settles every scientific
choice \parencite[pp.~747--753]{AndersenAndersen2021Scientific}.

The source uses \emph{science} for the natural sciences, in line with
the English-language discussions it surveys. Some of its review-policy
examples also concern mathematics, humanities, and social sciences.
Those examples do not, without further argument, extend every
natural-science account to every discipline. The historical and
philosophical positions below remain attributed positions; no
population-level estimate of scientific rationality is supplied.

\section{Two objectivity questions and two contexts}
\label{sec:p11-rat141-individual}

\begin{hbdefinition}{def:p11-rat141-objectivity}{Objectivity of theories and objectivity of methods}
The source distinguishes two ways of calling science objective.
\emph{Theory objectivity} concerns whether theories faithfully describe
facts about the world. \emph{Method objectivity} concerns whether the
methods used in inquiry are independent of personal and social values
and biases.

These are different assessments. A biased process can happen to reach
a true conclusion; an appropriately controlled inquiry can still yield
an erroneous conclusion. Identifying the two senses requires an
argument about the method's connection to the world, not merely the
same adjective. The chapter's later community accounts also challenge
the proposed independence-from-values interpretation of method
objectivity, rather than quietly adopting it as a universal definition.
\end{hbdefinition}

The authors characterize much early twentieth-century philosophy of
science as treating rational adherence to objective methods as producing
objective theories. The central question was consequently how an
individual should validate a scientific belief. This is their account
of a methodological ideal, not a claim that every historical scientist
followed it or that its success has been experimentally established
\parencite[p.~747]{AndersenAndersen2021Scientific}.

\begin{hbdefinition}{def:p11-rat141-discovery-justification}{Discovery and justification in the surveyed individualist ideal}
The \emph{context of discovery} concerns how a scientist forms a new
hypothesis. On the surveyed distinction, this is an empirical subject
for psychology, sociology, and history. The \emph{context of
justification} concerns how hypotheses should be assessed and how a
scientist should choose among competing theories. On that account,
this is a normative philosophical question.

The division separates a causal story of arriving at a claim from its
justification. It does not itself show that the actual discovery and
assessment activities never interact. The source's note~1 directs to
Hoyningen-Huene for different interpretations of the distinction.
\end{hbdefinition}

If a special scientific method fixes the rational assessment, adherence
to it supplies the proposed individual-rationality standard. The chapter
selects Hempel's hypothetico-deductive account and Popper's
falsificationism as influential versions. Bayesianism is named but
explicitly not examined in the selected extent; the source directs to
its chapter~4.1 instead. A probability update is therefore not inserted
as though it were the argument this chapter chose to develop.

\section{Confirmation, refutation and challenges to a unique method}
\label{sec:p11-rat141-values-and-methods}

\subsection{The hypothetico-deductive and falsificationist contrast}

In the source's account of Hempel's hypothetico-deductive, or H--D,
model, the scientist derives observable predictions from a hypothesis
and tests whether those events occur. Observed predictions provide
support whose degree depends on the amount, variety, and precision of
the evidence. Failed predictions count against the hypothesis. This
is a logic of confirmation, not the deduction of the hypothesis from
every successful observation
\parencite[pp.~747--748]{AndersenAndersen2021Scientific}.

Popper's surveyed alternative emphasizes refutation. True predicted
consequences do not, by their logical form alone, establish an empirical
hypothesis with certainty; a false consequence of a true premise is
impossible in a valid deduction. Rational inquiry should therefore
seek potential falsifiers and attempt sincerely to falsify its
conjectures. A hypothesis that makes more vulnerable empirical claims
offers more opportunities for refutation. Surviving tests is not the
same as verification. On the account presented here, survival motivates
refinement and more stringent tests rather than a claim that the
hypothesis has been conclusively established.

\begin{hbtheorem}{thm:p11-rat141-test-logic}{proved=annex:p11.rat141.test-logic}{What successful and failed predictions entail}
Let \(H\), \(A\), and \(E\) be independently assignable
propositional atoms in classical Boolean semantics, representing a
hypothesis, the conjunction of test auxiliaries, and a predicted
observation. Then
\[
 (H\land A)\to E,\ \neg E\ \entails\ \neg(H\land A),
 \qquad
 (H\land A)\to E,\ A,\ \neg E\ \entails\ \neg H.
\]
The first two premises alone do not entail \(\neg H\).
Nor do \((H\land A)\to E\), \(A\), and \(E\) entail \(H\).
Thus a successful consequence does not deductively verify the
hypothesis, and a failed consequence refutes the tested conjunction;
isolating the hypothesis as the rejected premise additionally requires
the auxiliaries.
\end{hbtheorem}

This is an explicit logical boundary of the source's prediction/refutation
contrast, not a claim that Hempel or Popper ignored every auxiliary
assumption. It matters in an actual test: calibration, initial
conditions, sample preparation, and the description of the observation
can enter the tested premise set. A failure then presents a question
about which premise to revise, not a deductive instruction to discard
only the favored hypothesis. Likewise, if a statistical hypothesis
merely gives an event a small probability, its occurrence is not a false
deductive consequence. A tail-probability assessment needs a specified
sampling law and inferential standard.

The logical point also has limits. The countervaluation proves that
the indicated inference form is invalid in general; it does not prove
that no hypothesis can ever be deduced from sufficiently informative
premises. The source's anti-verification point concerns scientific
empirical hypotheses and their successful tests, not tautologies or a
premise set that already entails the hypothesis independently.

The source reports the influence of unique-method ideals beyond
philosophy, including textbook presentations that move from a
phenomenon, to a proposed explanation, to an experimental test. It
does not establish that this sequence exhausts scientific practice.
Its note~2 gives Potochnik, Colombo, and Wright's textbook as a
counterexample to the impression that all teaching presents one
method. Its references to scientific-method histories and surveys
identify further discussions rather than import their entire accounts
\parencite[p.~748 and p.~753, nn.~1--2]{AndersenAndersen2021Scientific}.

\subsection{Values, historical variation and theory-laden observation}

The first challenge concerns accepting a hypothesis as a decision
with possible error costs. Churchman and Rudner, as surveyed by the
authors, require attention both to evidential strength and to the
consequences of being mistaken in acceptance. The chapter thereby
distinguishes an epistemic question---which theory is warranted by
the evidence---from an instrumental question---what to do given
beliefs and ends. A practical consequence does not automatically change
the evidence, but can matter to the decision about how to proceed.
Whether epistemic alone or also social, ethical, and political values
may legitimately influence theory choice is itself disputed.

The authors connect this with the Fisher/Neyman--Pearson controversy.
They characterize Fisher's position as seeking a numerical expression
of disbelief and Neyman and Pearson's positions as including error
consequences. This is an attributed historical characterization, not
a definition that equates a modern significance-test tail probability
with the probability that a hypothesis is false. Their selected survey
does not reproduce those original statistical methods or settle the
debate. It points to the value-freedom discussion in the source's
chapter~14.2 for further treatment
\parencite[p.~748]{AndersenAndersen2021Scientific}.

The second challenge is historical and disciplinary variation in aims
and norms. The source names Feyerabend, Hanson, Kuhn, Laudan, Shapere,
and Toulmin as historically oriented critics of a universal,
ahistorical rational method. Feyerabend's methodological anarchism is
the radical version: a methodological rule can sometimes be fruitfully
violated. This is a claim about the limits of universal rules, not
permission to infer that every procedure is equally good or that
arbitrary error is productive. Note~3 names fuller histories and
accounts of these positions.

The third challenge is theory-laden observation. An observation can
depend on theoretical assumptions, so the evidential relation is more
complex than an uninterpreted fact simply checking an isolated theory.
The point does not establish that there is no world constraining the
inquiry. It requires making the interpretation and background
assumptions visible. Whether they are justified remains a question
that cannot be answered merely by observing that they are present
\parencite[pp.~748--749 and p.~753, n.~3]{AndersenAndersen2021Scientific}.

\section{The community as the unit of scientific rationality}
\label{sec:p11-rat141-community-rationality}

\subsection{Kuhn's essential tension and rational disagreement}

The authors present Kuhn as shifting the central evaluand from each
individual to the scientific community. Normal science works within
a received paradigm that specifies legitimate problems and acceptable
approaches and solutions. Recalcitrant anomalies can produce crisis
and eventually a revolution. The new paradigm need not cumulatively
extend the old: the earlier framework may have excluded the very kind
of solution the successor supplies. Incommensurability therefore blocks
a straightforward common comparison, not every possible discussion
of the alternatives.

The source's explanation of why this need not be irrational appeals
to the essential tension. Stable normal-science practice concentrates
effort on defined problems; fundamental innovations also require
challenges to that practice. A distribution of conservative and
innovative dispositions permits some researchers to pursue alternatives
while others continue established work. Note~4 points to Hanne
Andersen's fuller discussion of the mechanism. This is a philosophical
account of the balance; no calculated optimal diversity level or
measured transition time is established in the selected chapter
\parencite[p.~749]{AndersenAndersen2021Scientific}.

\begin{hbdefinition}{def:p11-rat141-kuhn-values}{Shared values that constrain without uniquely determining theory choice}
The Kuhnian values identified by the source are accuracy, consistency,
scope, simplicity, and fruitfulness. Members of a scientific community
can share these values while interpreting them or giving them different
weights. Values can conflict. Sharing the list does not therefore
determine one decision by all rational individuals.

The source also allows personal preferences concerning originality,
risk taking, aesthetics, and philosophical commitments to affect which
theory a scientist pursues. It does not assign numerical weights or
provide an aggregation formula for these considerations.
\end{hbdefinition}

\begin{hbexample}{def:p11-rat141-value-conflict}{The source's simplicity--accuracy and originality contrasts}
The simpler of two competing theories need not be the more accurate.
Researchers who agree that both simplicity and accuracy matter can
therefore disagree if they weigh or interpret those criteria differently.
This differs from one researcher overlooking a relevant datum or
making an invalid deduction; the example supposes a legitimate
evaluative conflict.

The source's later review discussion adds Kuhn's example of originality:
a scientist valuing it more may be more prepared to take epistemic
risks and may reasonably use a lower evidential threshold for accepting
a hypothesis. Whether that risk is reasonable depends on the activity
and its consequences. The example is not a rule licensing a lower
threshold for every practical application of a novel claim.
\end{hbexample}

On the surveyed model, an alternative capable of resolving anomalies
can gradually win support until the community reaches a new consensus.
The consensus reflects the proliferation of arguments as scientists
assess the alternative through their own interpretations of the shared
values. It is the community's effective decision, not the execution of
one uniquely determining individual algorithm. The authors contrast
this with the earlier unique-method ideal
\parencite[pp.~749--750 and p.~753]{AndersenAndersen2021Scientific}.

Consensus nevertheless is not made identical to truth or objectivity.
The source acknowledges that Kuhn at most hints at a corresponding
community account of objectivity: accuracy and consistency should
figure in it, although they supply less determinate guidance than the
traditional picture expected. His detailed rationality narrative and
this less-developed objectivity suggestion have different explanatory
standing. The chapter reports objections by Lakatos, Scheffler, and
Shapere that the picture makes theory choice look irrational. Its
response is the argument about shared values and community development,
not a proof that disagreement or eventual unanimity alone tracks truth.

The strong program in sociology of science, associated here with Bloor,
is a different claim. It proposes explaining successful and unsuccessful
knowledge claims impartially through the same kinds of social causes.
Its explanatory symmetry is not the claim that the beliefs have equal
evidential warrant. Nor does the presence of social causes by itself
remove the distinction between a cause of believing and a reason for
the belief \parencite[p.~750]{AndersenAndersen2021Scientific}.

\subsection{Credit, incentives and the division of inquiry}

Hull's surveyed account links the individual desire for credit to
the community's epistemic goals. Reward for true contributions and
severe punishment for detected fraud are meant to make pursuit of
truth self-interested. The explanation depends on the actual reward,
detection, and sanction structure; a desire for credit by itself would
not prevent fraud. The source's strong formulation that this structure
ensures scientists will not cheat is therefore an account of the
proposed incentive mechanism, not a verified claim of zero scientific
misconduct \parencite[p.~750]{AndersenAndersen2021Scientific}.

\begin{hbexample}{def:p11-rat141-unorthodox-method}{Kitcher's fame-motivated choice of a different method}
Suppose researchers can pursue a problem using several methods. All
may be tempted to use the method judged intrinsically best. Yet the
solution might be found by an unorthodox method. A researcher seeking
fame can have a reason to pursue that method: it retains some chance
of success and has fewer direct competitors.

The source uses this case to explain how a nonepistemic individual
motive can support an epistemically useful division of scientific
effort. The conclusion requires the alternative method to have a
relevant chance of success and the incentive to favor genuinely useful
exploration. Fame can also reward publicity or unreliable claims;
the mere motive does not guarantee the favorable outcome.
\end{hbexample}

The chapter names game-theoretic, economic, and mathematical-modeling
work on heterogeneous communities by D'Agostino, De Langhe,
Weisberg, Zollman, and collaborators. It explicitly includes the
possibility that heterogeneity is beneficial \emph{or detrimental}.
These references do not supply one theorem that any diversity always
improves science. No unspecified game, learning rule, landscape, or
network result is treated as locally proved. The source refers to
social epistemology for the general community-level framework
\parencite[pp.~750--751]{AndersenAndersen2021Scientific}.

\section{Group objectivity and the conditions of effective criticism}
\label{sec:p11-rat141-community-objectivity}

The source's feminist discussion starts with the situated and
perspectival character of knowledge: what a person knows reflects
their position; there is no neutral view detached from every knower.
The authors describe criticism of traditional objectivity models that
disadvantage women and minority groups. Responses differ. The
surveyed Irigaray position embraces subjective ways of knowing nature;
the Longino, Haraway, and Solomon approaches develop community-level
objectivity instead. These are not one conclusion forced by the word
\emph{situated} \parencite[p.~751]{AndersenAndersen2021Scientific}.

\begin{hbdefinition}{def:p11-rat141-contextual-evidence}{Background assumptions and contextual evidential support}
In the chapter's reconstruction of Longino, data bear on a theory
relative to background assumptions. The theory states more than the
data description alone, so the data do not determine all its claims.
The assumptions mediate the evidential connection and can carry
subjective biases.

This is an account of epistemic justification, not the assertion that
any theory is supported by any data. It makes the background premises
available for criticism rather than declaring that their presence
automatically invalidates every inference.
\end{hbdefinition}

Longino's question is how inquiry can be objective when individuals'
background assumptions are biased. Her proposed answer is that results
must be evaluated by a plurality of individuals and shaped by their
critical exchanges. This is not merely adding up completed individual
products. Other researchers repeat experiments with variations,
criticize and reformulate hypotheses, and modify proposed results
before they enter accepted knowledge. The mechanism is the exposure,
challenge, and revision of assumptions that any one perspective may
leave unseen \parencite[p.~751]{AndersenAndersen2021Scientific}.

\begin{hbdefinition}{def:p11-rat141-critical-process}{The source's explicit community-objectivity criteria}
On the Longino-style account selected here, objectivity admits degrees
and depends on a community's capacity to elicit multiple perspectives
in the critical discussion that forms its results. The selected chapter
explicitly identifies two requirements from Longino's tentative list:
appropriate venues for criticism and a community that responds to
criticism.

The requirements concern effective participation and uptake, not
headcount alone. Presenting several views while ignoring their objections
does not implement the described process. The source cites Longino
2006, pp.~172--173, for the fuller list; it does not enumerate that
entire list here. These two named requirements are therefore not
presented as a complete independently proved sufficient criterion.
\end{hbdefinition}

The source reports Longino as holding existing scientific inquiry
generally objective while also recommending better organization to
increase its objectivity. That evaluative claim is not a survey
estimate, nor does it guarantee that every journal or laboratory
satisfies the critical-process conditions. Diversity must actually
enter the reasoning and affect responses. The passage's claim that
more perspectives increase objectivity belongs to this effective
critical-process account; it is not a measured monotonic relation
between demographic counts and accuracy.

Haraway's surveyed situated-knowledges account emphasizes embodied
and partial perspectives, rejecting the claim of an unconstrained
observer who sees from nowhere. Harding's strong-objectivity proposal
adds the claim that some marginalized perspectives have an epistemic
advantage. The latter is an additional substantive thesis; it does
not follow just from everybody having a perspective. The chapter
also names Solomon and Wylie as related contributors without
reproducing each separate theory
\parencite[p.~751]{AndersenAndersen2021Scientific}.

Note~5 identifies Longino's human-evolution and behavioral-
neuroendocrinology cases and Lee and colleagues' review of peer-review
bias. They are examples and evidence sources \emph{cited by the
chapter}, not newly acquired participant studies here. No result
with an unstated sample or measurement is inferred from these titles.
Note~6 preserves two useful affinities without making the accounts
identical: Popper's intersubjective-testability characterization, and
Habermas's justification of truth claims through critical discussion
among free and uncoerced participants whose reasons could be accepted
by any reasonable person. Shared emphasis on criticism does not erase
their different accounts of evidence, background assumptions, or
justification \parencite[p.~753, nn.~5--6]{AndersenAndersen2021Scientific}.

On a Longino-style account, peer review can introduce alternative
perspectives as well as check data and arguments. Postpublication
replication, modification, and integration with other results can also
shape what is eventually accepted. The rationality of this process
therefore cannot be assessed solely at the instant an individual first
submits a claim.

\section{Peer review as the chapter's practical test case}
\label{sec:p11-rat141-peer-review}

\begin{hbdefinition}{def:p11-rat141-peer-review}{Single-blind, double-blind and open review in the selected source}
In \emph{single-blind} review the reviewer knows the author's identity
but the author does not know the reviewer's. In \emph{double-blind}
review both identities are concealed from the other party. In
\emph{open} review the parties know one another's identities; some
systems additionally publish reviewer names, reviews, or other parts
of the exchange.

Knowing the parties' identities and publishing the exchange are
distinct design choices. The labels do not themselves establish which
information is accessible to the wider community or whether anonymity
is effective.
\end{hbdefinition}

The source reports single-blind review as common and double-blind
review as especially common in humanities and social sciences, using
Lee and colleagues and Ware as its references. These are dated
literature reports, not a newly measured current prevalence. The
chapter supplies no sampling frame, task-specific effect estimate,
or uncertainty interval for those general prevalence claims
\parencite[p.~752]{AndersenAndersen2021Scientific}.

Concealing reviewers' identities is defended as protection against
repercussions and a condition for frank criticism. Concealing authors'
identities is defended as protection against reviewers' social biases.
On the individual-method objectivity picture, anonymity can keep
irrelevant social information from distorting method-based assessment.
If reviewers were already perfectly impartial and safe from retaliation,
these particular reasons for blinding would lose their point. Thus
the case depends on actual risks, not anonymity having an intrinsic
truth-making property.

From the community-objectivity picture, concealment can hide the
theoretical and cultural positions that help participants detect one
another's background assumptions. Not knowing the author can make
it harder for a reviewer to identify assumptions in the manuscript;
not knowing the reviewer can similarly hinder the author's criticism
of the review. Disclosure may therefore help the critical exchange.
Yet it can also suppress honest criticism when a reviewer fears
retaliation. The authors expressly decline to resolve this tradeoff.
Their conclusion is that the account of objectivity changes the
considerations relevant to review design, not that open review always
wins \parencite[p.~752]{AndersenAndersen2021Scientific}.

\begin{hbexample}{def:p11-rat141-bmj}{The source's public-review example and its limits}
The chapter describes a BMJ practice, citing Groves and Loder 2014
and Smith 1999/2006, that publishes reviewer identities and comments,
author responses, and article versions online. The intended benefits
are transparency and accountability. Readers can examine the review
as a perspective in public critical discussion, and public visibility
can give reviewers an incentive to be thorough and fair.

This is the source's dated institutional example, not a verified
description of every present BMJ article or a measured effect of
disclosure. Assessing the proposed benefit would also require checking
participation, candor, uptake, retaliation, and what actually becomes
visible. Public availability alone cannot demonstrate effective criticism.
\end{hbexample}

Note~7 adds an important limit on double blindness: a specialized
manuscript's content can disclose its author. Nominal concealment and
effective concealment are different properties. The source points to
mathematics and chemistry discussions and further references, without
supplying a frequency estimate of successful identification
\parencite[p.~752 and p.~753, n.~7]{AndersenAndersen2021Scientific}.

\subsection{Agreement is not yet reliability}

The source challenges the assumption that agreement between reviewers
is always desirable. Lee's Kuhnian argument allows appropriate
differences in interpreting and applying soundness, significance, and
novelty; the originality/risk case is its worked illustration. On a
Longino-style account, high agreement can instead reflect common
background biases. Low agreement can be compatible with a community
that elicits effective criticism. These possibilities show why an
agreement statistic needs an interpretation rather than providing a
universal objectivity score.

The distinctions run in both directions. Disagreement can also arise
from mistakes, poor attention, or unequal information. Agreement can
reflect sound examination of the same evidence. Neither observation
alone identifies its cause. A false proof is not rendered valid by
diverse evaluative priorities, and a well-supported result is not
discredited just because several informed reviewers agree.

The chapter reports that empirical studies have found low
inter-reviewer agreement, but does not give an individual study's
sample, review task, agreement measure, estimate, uncertainty, or
replication details. The warranted local conclusion is consequently
the philosophical one: low agreement is \emph{not necessarily} a
reason for concern on every account of objectivity. No numerical
agreement finding, general effect of open review, or diagnosis of a
particular journal is asserted here. An empirical investigation would
need the evidence record of Method~\ref{def:p19-evidence-record}
\parencite[pp.~752--753]{AndersenAndersen2021Scientific}.

\section{Using the standards and keeping the handbook interfaces distinct}
\label{sec:p11-rat141-use}

\begin{hbmethod}{def:p11-rat141-objectivity-check}{Specify the agent, the standard and the mechanism}
First state what is being assessed: a hypothesis's truth or warrant,
an individual's test method, a research choice, a community's division
of inquiry, a critical exchange, or a review policy. Distinguish
theory objectivity from the proposed process criterion.

State the actual standard. For an H--D test identify the prediction,
test auxiliaries, and evidence; distinguish support from entailment.
For falsification identify the excluded observation and which premises
its occurrence would refute. For Kuhnian inquiry specify the shared
values and their disputed interpretation or weighting. For community
objectivity identify the perspectives, the criticism venues, and
the actual responses to criticism.

Give the mechanism rather than a favorable label. If incentives are
said to favor truth, explain credit, detection, sanctions, and the
available alternatives. If anonymity is said to reduce bias or openness
to expose assumptions, specify what information is concealed or
revealed, to whom, and how it affects criticism. Account separately
for retaliation and effective concealment.

Keep normative premises distinct from descriptive findings. Agreement,
diversity, transparency, and eventual acceptance are observations or
institutional properties; none by itself establishes truth or a good
scientific decision. If an empirical policy effect is asserted, supply
its task, population, design, measure, uncertainty and alternatives,
rather than transferring an example's rationale to every discipline.
\end{hbmethod}

This is an editorial application of the source's distinctions. Its
success case is a peer-review assessment that makes both its critical-
exchange objective and its candor/retaliation constraints explicit,
then examines whether criticisms can actually change the result.
Its failure cases include calling a consensus objective solely because
it is unanimous, calling disagreement rational solely because it is
diverse, and rejecting a hypothesis from a failed prediction while
silently leaving untested auxiliaries unquestioned. These concern
different standards and do not share one automatic repair.

The H--D account in this source is not the same question as Hempel
and Oppenheim's 1948 deductive-nomological explanation in
Definition~\ref{def:p16-dn-schema}. The latter requires an explanandum,
laws and conditions, empirical content and true explanatory premises;
its valid derivation does not determine the warrant for those premises
or the choice between competing scientific theories. Explanation,
prediction, and hypothesis confirmation have connected but distinct
tasks.

P15's Definition~\ref{def:p15-finite-hypothesis-partition} separates
likelihood of evidence under a hypothesis from the posterior weight
of the hypothesis, and Theorem~\ref{thm:p15-finite-bayes} gives the
exact finite update with its positivity assumptions. P17's
Definition~\ref{def:p17-binary-test} supplies two specified error
probabilities. Neither object is, by itself, a complete measure of
scientific objectivity or a justification of the costs assigned to
errors. The practical/theoretical distinction in
Definition~\ref{def:p11-rat22-attitude-types} keeps evidence-supported
belief apart from a decision about what to do.

For claims about how review policies \emph{cause} changes in performance,
the observational/intervention distinction in
Chapter~\ref{sec:p16-finite-causal-models} matters. An observational
association between openness and acceptance need not identify the
effect of changing a policy. Defining an intervention model also does
not establish that its assumptions describe the institution under study.

Social epistemology's
Definition~\ref{def:p18-r101-holders} distinguishes knowledge in
groups from aggregate, common, and corporate group attitudes.
Its Definition~\ref{def:p18-r101-correctness} requires a correctness
standard independent of the procedure being evaluated. The present
chapter's community inquiry and transformative criticism cannot be
substituted for a theorem about majority accuracy, nor can a majority
theorem certify its sociological or normative premises. P19's
Definition~\ref{def:p19-benchmark} similarly separates a benchmark
from actual observed responses.

The source concludes with a contrast among accounts, not one admitted
policy: some place rationality and objectivity in an individual's
adherence to method; others explain them through differently motivated
individuals contributing to a community's inquiry and criticism. The
questions of method, values, incentives, plural perspectives and
institutional uptake remain connected without becoming identical.
The only local theorem proves the stated logical boundary; the
remaining arguments have their historical, philosophical and
normative standing \parencite[p.~753]{AndersenAndersen2021Scientific}.

\index{scientific rationality}
\index{scientific objectivity!individual and community}
\index{hypothetico-deductive method}
\index{falsification}
\index{scientific disagreement}
\index{peer review!anonymity and openness}
